\documentclass{article}
\begin{document}
\section*{Binary native carrier-changing naturality square}
The source and target each have a fixed ordered occurrence-first
and joint finite-family presentation. The comparison reindexes the
identical raw tags and applies the invertible multiplicative middle-four
on each resource. The two horizontal arrows are independently built
from the supplied whole parent maps and dependent continuation maps:
no individual source resource map is chosen. These are actual native
arrows for the standard symmetric additive direct image.

A source tuple lies over a target tuple precisely when each of its
parent and parent-dependent continuation coordinates maps to the
chosen one. The selection is proved as exact list equality in both
orders. A local continuation is demanded only from a source parent
in the selected preimage, and only for a valid target continuation.
Stable filtering retains order; distinct raw tags can have equal
coefficient values. A local aggregate from the additive unit is still
a supplied actual arrow. A concrete many-to-one example verifies
that two tags survive one fibre and another fibre is genuinely empty.

The native equality is reduced to the restricted tagged reorder and
the sum of the actual leaf shuffles, followed by the two opaque
whole-fibre routes. The fixed-fibre cell equation proves this equality.
The assumptions are distinct raw position tags, valid support maps,
the declared binary rig comparison laws, and invertibility of the
multiplicative middle-four. This is naturality as a square of native
morphisms, not an assertion that the horizontal constructions are
already functors. Preservation of global augmentations by the
horizontal arrows, their identity and composition laws, and the
all-arity dependent theorem are separate obligations.
\end{document}
